\documentclass[12pt]{article}
\usepackage[UTF8]{inputenc}
\usepackage{thaisupp}
\usepackage{enumerate}
\usepackage{etoolbox}
\AtBeginEnvironment{enumerate}{\itshape}
\begin{document}
\title{การเคลื่อนที่แบบวงกลม}
\maketitle
\section*{คำถาม in class}
\begin{enumerate}
\def\labelenumi{\arabic*.}
\normalfont
\item \textbf{วัตถุมวล 2 kg เคลื่อนที่เป็นวงกลมรัศมี 0.5 m ด้วยความเร็ว 4 m/s แรงสู่ศูนย์กลางมีขนาดเท่าใด}
\begin{enumerate}[label=\alph*)]
\item 8 N
\item 16 N
\item 32 N
\item 64 N
\end{enumerate}
\item \textbf{วัตถุเคลื่อนที่เป็นวงกลมครบ 10 รอบในเวลา 20 วินาที ความเร็วเชิงมุมเป็นเท่าใด}
\begin{enumerate}[label=\alph*)]
\item 0.5 rad/s
\item 1 rad/s
\item 2 rad/s
\item π rad/s
\end{enumerate}
\item \textbf{ดาวเทียมโคจรรอบโลก 1 รอบใช้เวลา 2 ชั่วโมง คาบการโคจรเป็นเท่าใด}
\begin{enumerate}[label=\alph*)]
\item 2 ชั่วโมง
\item 4 ชั่วโมง
\item 7,200 วินาที
\item ทั้ง ก และ ค
\end{enumerate}
\item \textbf{จานหมุนหมุนได้ 120 รอบต่อนาที ความถี่ของการหมุนเป็นเท่าใด}
\begin{enumerate}[label=\alph*)]
\item 0.5 Hz
\item 2 Hz
\item 60 Hz
\end{enumerate}
\item \textbf{วัตถุมวล 1 kg ผูกเชือกยาว 0.5 m หมุนเป็นวงกลมในระนาบดิ่งด้วยความเร็ว 3 m/s ที่จุดต่ำสุด แรงตึงเชือกเป็นเท่าใด (g = 10 m/s²)}
\begin{enumerate}[label=\alph*)]
\item 9 N
\item 10 N
\item 18 N
\item 28 N
\end{enumerate}
\item \textbf{รถยนต์มวล 1,000 kg วิ่งเข้าโค้งรัศมี 100 m ด้วยความเร็ว 20 m/s ต้องมีแรงเสียดทานอย่างน้อยเท่าใดจึงจะไม่ไถล (g = 10 m/s²)}
\begin{enumerate}[label=\alph*)]
\item 1,000 N
\item 2,000 N
\item 4,000 N
\item 8,000 N
\end{enumerate}
\item \textbf{วัตถุเคลื่อนที่เป็นวงกลมรัศมี 2 m ด้วยความเร็วเชิงมุม 4 rad/s ความเร็วเชิงเส้นเป็นเท่าใด}
\begin{enumerate}[label=\alph*)]
\item 2 m/s
\item 4 m/s
\item 8 m/s
\item 16 m/s
\end{enumerate}
\item \textbf{เครื่องปั่นหมุน 3,000 รอบต่อนาที ความถี่เป็นเท่าใด และคาบเป็นเท่าใด}
\begin{enumerate}[label=\alph*)]
\item 30 Hz, 0.033 s
\item 50 Hz, 0.02 s
\item 60 Hz, 0.0167 s
\item 300 Hz, 0.0033 s
\end{enumerate}
\item \textbf{วัตถุมวล 0.5 kg ผูกเชือกยาว 1 m หมุนเป็นวงกลมในระนาบราบด้วยความเร็ว 5 m/s แรงตึงเชือกเป็นเท่าใด}
\begin{enumerate}[label=\alph*)]
\item 5 N
\item 10 N
\item 12.5 N
\item 25 N
\end{enumerate}
\item \textbf{อิเล็กตรอนเคลื่อนที่เป็นวงกลมในสนามแม่เหล็ก ถ้ารัศมีวงโคจรเพิ่มขึ้น 2 เท่า โดยความเร็วคงที่ แรงสู่ศูนย์กลางจะเป็นเท่าใด}
\begin{enumerate}[label=\alph*)]
\item เพิ่ม 2 เท่า
\item ลดลง 2 เท่า
\item คงที่
\item เพิ่ม 4 เท่า
\end{enumerate}
\item \textbf{ล้อรัศมี 0.3 m หมุนด้วยความเร็วเชิงเส้นที่ขอบ 6 m/s ความเร็วเชิงมุมเป็นเท่าใด}
\begin{enumerate}[label=\alph*)]
\item 2 rad/s
\item 5 rad/s
\item 10 rad/s
\item 20 rad/s
\end{enumerate}
\item \textbf{วัตถุมวล 2 kg ผูกเชือกยาว 0.8 m หมุนเป็นวงกลมในระนาบดิ่ง ที่จุดสูงสุดด้วยความเร็ว 6 m/s แรงตึงเชือกเป็นเท่าใด (g = 10 m/s²)}
\begin{enumerate}[label=\alph*)]
\item 30 N
\item 50 N
\item 70 N
\item 90 N
\end{enumerate}
\item \textbf{รถแข่งวิ่งบนไบร์ตรัศมี 50 m ต้องมีความเร็วอย่างน้อยเท่าใดจึงจะไม่ตกออกนอกไบร์ต (g = 10 m/s²)}
\begin{enumerate}[label=\alph*)]
\item 10 m/s
\item 15 m/s
\item 20 m/s
\item 22.4 m/s
\end{enumerate}
\item \textbf{วัตถุมวล 1 kg ผูกเชือกยาว 0.5 m หมุนเป็นวงกลมในระนาบดิ่ง ถ้าแรงตึงเชือกสูงสุดที่ทนได้คือ 50 N วัตถุจะหมุนได้โดยไม่ขาดที่ความเร็วสูงสุดเท่าใด (g = 10 m/s²)}
\begin{enumerate}[label=\alph*)]
\item 4.5 m/s
\item 5.0 m/s
\item 5.5 m/s
\item 6.0 m/s
\end{enumerate}
\item \textbf{มอเตอร์หมุนที่ความถี่ 1,800 รอบ/นาที ความเร็วเชิงมุมเป็นเท่าใด}
\begin{enumerate}[label=\alph*)]
\item 30π rad/s
\item 60π rad/s
\item 90π rad/s
\item 120π rad/s
\end{enumerate}
\item \textbf{ดาวเทียมมวล 500 kg โคจรรอบโลกที่ความสูงเท่ากับรัศมีโลก (จากจุดศูนย์กลางโลก = 2R) ด้วยความเร็ว 5,600 m/s จงหาแรงสู่ศูนย์กลาง (R = 6.4×10⁶ m)}
\begin{enumerate}[label=\alph*)]
\item 1,225 N
\item 2,450 N
\item 4,900 N
\item 9,800 N
\end{enumerate}
\item \textbf{ดาวเทียมโคจรรอบโลกที่ความสูง 3R จากผิวโลก (R = รัศมีโลก) ถ้าคาบของดาวเทียมใกล้พื้นผิวโลก = 84 นาที คาบของดาวเทียมนี้เป็นเท่าใด}
\begin{enumerate}[label=\alph*)]
\item 84 นาที
\item 168 นาที
\item 336 นาที
\item 672 นาที
\end{enumerate}
\item \textbf{วัตถุมวล 0.5 kg ผูกเชือกยาว 0.8 m หมุนเป็นวงกลมในระนาบดิ่ง ที่จุดกึ่งกลาง (ระดับเดียวกับจุดศูนย์กลาง) ด้วยความเร็ว 6 m/s แรงตึงเชือกเป็นเท่าใด (g = 10 m/s²)}
\begin{enumerate}[label=\alph*)]
\item 22.5 N
\item 27.5 N
\item 45 N
\item 50 N
\end{enumerate}
\item \textbf{อนุภาคมวล 0.1 kg เคลื่อนที่เป็นวงกลมรัศมี 0.5 m ด้วยความเร็วเชิงมุม 10 rad/s แรงสู่ศูนย์กลางเป็นเท่าใด}
\begin{enumerate}[label=\alph*)]
\item 0.5 N
\item 2.5 N
\item 5 N
\item 50 N
\end{enumerate}
\item \textbf{วัตถุเคลื่อนที่เป็นวงกลมครบ 1 รอบในเวลา 0.5 วินาที ความถี่เป็นเท่าใด}
\begin{enumerate}[label=\alph*)]
\item 0.5 Hz
\item 1 Hz
\item 2 Hz
\item 5 Hz
\end{enumerate}
\item \textbf{รถบรรทุกมวล 4,000 kg วิ่งเข้าโค้งรัศมี 200 m ด้วยความเร็ว 72 km/hr ต้องการแรงเสียดทานเป็นเท่าใดจึงจะไม่ไถล (g = 10 m/s²)}
\begin{enumerate}[label=\alph*)]
\item 4,000 N
\item 8,000 N
\item 16,000 N
\item 32,000 N
\end{enumerate}
\item \textbf{เข็มนาฬิกายาว 1 cm ปลายเข็มเคลื่อนที่เป็นวงกลมรัศมี 1 cm ครบ 1 รอบใน 60 วินาที ความเร็วเชิงเส้นที่ปลายเข็มเป็นเท่าใด}
\begin{enumerate}[label=\alph*)]
\item π/30 m/s
\item π/60 m/s
\item π/15 m/s
\item π/3 m/s
\end{enumerate}
\item \textbf{วัตถุมวล 3 kg ผูกเชือกยาว 1.5 m หมุนเป็นวงกลมในระนาบดิ่ง ที่จุดสูงสุดมีแรงตึงเชือก 10 N ความเร็วที่จุดสูงสุดเป็นเท่าใด (g = 10 m/s²)}
\begin{enumerate}[label=\alph*)]
\item 3.5 m/s
\item 5.0 m/s
\item 7.0 m/s
\item 10.0 m/s
\end{enumerate}
\item \textbf{ลูกบอลมวล 0.2 kg ผูกเชือกยาว 0.4 m แกว่งเป็นวงกลมในระนาบราบ 2 รอบต่อวินาที แรงตึงเชือกเป็นเท่าใด}
\begin{enumerate}[label=\alph*)]
\item 4 N
\item 8 N
\item 16 N
\item 32 N
\end{enumerate}
\item \textbf{จานเสียงหมุน 33 รอบต่อนาที ถ้าต้องการให้จานหมุนเร็วขึ้นเป็น 45 รอบต่อนาที โดยไม่เปลี่ยนความเร็วเชิงเส้นที่ขอบ ต้องเปลี่ยนรัศมีเป็นกี่เท่าของรัศมีเดิม}
\begin{enumerate}[label=\alph*)]
\item 0.73 เท่า
\item 1.37 เท่า
\item 1.5 เท่า
\item 2 เท่า
\end{enumerate}
\end{enumerate}
\end{document}